% !TEX program = xelatex
\documentclass{英语试卷}
\begin{document}
\examtitle{2025年3月济南市高三第一次模拟考试}{英语}
\begin{notice}
\par 1. 答卷前，考生务必用黑色字迹的钢笔或签字笔将自己的姓名、考生号、试室号和座位号填写在答题卡上。用 2B 铅笔将试卷类型（A）填涂在答题卡相应位置上，并在答题卡相应位置上填涂考生号。\par
\par 2. 作答选择题时，选出每小题答案后，用 2B 铅笔把答题卡对应题目选项的答案信息点除黑；如需改动，用橡皮擦干净后，再选涂其他答案。答案不能答在试卷上。\par
\par 3. 非选择题必须用黑色字迹的钢笔或签字笔作答，答案必须写在答题卡各题目指定区域内相应位置上，如需改动，先划掉原来的答案，然后再写上新答案；不准使用铅笔和涂改液。不按以上要求作答无效。\par
\par 4. 考生必须保持答题卡的整洁。考试结束后，将试卷和答题卡一并交回。\par
\end{notice}
\parttitle{第一部分 听力 (共两节，满分 30 分)}
\par 注意，听力部分答题时，请先将答案标在试卷上。听力部分结束前，你将有两分钟的时间将答案转涂到答题卡上。\par
\sectitle{第一节 (共 5 小题; 每小题 1.5 分, 满分 7.5 分)}
\par 听下面 5 段对话。每段对话后有一个小题，从题中所给的 A、B、C 三个选项中选出最佳选项。听完每段对话后，你都有 10 秒钟的时间来回答有关小题和阅读下一小题。每段对话仅读一遍。\par
\question{1. What probably is the time now?}
\begin{choices}(3)
  \item 10:20.
  \item 10:35.
  \item 10:50.
\end{choices}
\question{2. Why does the boy need to get up early today?}
\begin{choices}(3)
  \item To go for a run.
  \item To attend a lecture.
  \item To meet a professor.
\end{choices}
\question{3. What will the man do next?}
\begin{choices}(3)
  \item Find a restaurant.
  \item Cook a meal.
  \item Perform a magic trick.
\end{choices}
\question{4. What is the destination of the speakers?}
\begin{choices}(3)
  \item Florida.
  \item Boston.
  \item Tokyo.
\end{choices}
\question{5. What can we learn about Professor Smith's lectures?}
\begin{choices}(3)
  \item They're well attended.
  \item They're uninspiring.
  \item They're chaotic.
\end{choices}
\sectitle{第二节 (共 15 小题; 每小题 1.5 分, 满分 22.5 分)}
\par 听下面 5 段对话或独白。每段对话或独白后有几个小题，从题中所给的 A、B、C 三个选项中选出最佳选项。听每段对话或独白前，你将有 5 秒钟的时间阅读各个小题；听完后，各小题将给出 5 秒钟的作答时间。每段对话或独白读两遍。\par
\par 听第 6 段材料，回答第 6、7 题。\par
\question{6. Where does the conversation probably take place?}
\begin{choices}(3)
  \item On the bus.
  \item On the subway.
  \item In the gym.
\end{choices}
\question{7. What caused the team's failure according to the woman?}
\begin{choices}(3)
  \item They made many mistakes.
  \item The opponents were strong.
  \item The key player was absent.
\end{choices}
\par 听第 7 段材料，回答第 8 至 10 题。\par
\question{8. Who will join the speakers?}
\begin{choices}(3)
  \item David and Maria.
  \item Peter and Allie.
  \item Jennifer and David.
\end{choices}
\question{9. Where is Walter?}
\begin{choices}(3)
  \item In London.
  \item In Paris
  \item In New York.
\end{choices}
\question{10. What will Michael do during the holiday?}
\begin{choices}(3)
  \item Stay at home.
  \item Visit his relatives.
  \item Travel to the desert.
\end{choices}
\par 听第 8 段材料，回答第 11 至 13 题。\par
\question{11. What is the woman mainly talking about?}
\begin{choices}(3)
  \item Her trip to Brazil.
  \item Her sales experience.
  \item Her special colleague.
\end{choices}
\question{12. Where did the staff find Tucson?}
\begin{choices}(3)
  \item At the car lot.
  \item At the pet shelter.
  \item Near the boss's house.
\end{choices}
\question{13. What happened after Tucson joined the sales team?}
\begin{choices}(3)
  \item He lost his ID card soon.
  \item He was hit by a customer.
  \item He enjoyed great popularity.
\end{choices}
\par 听第 9 段材料,回答第 14 至 17 题。\par
\question{14. What does the man seem to like most about the accommodation?}
\begin{choices}(3)
  \item The neighbors are friendly
  \item The transport is convenient
  \item The security system is good.
\end{choices}
\question{15. What color would the woman like her bedroom walls to be?}
\begin{choices}(3)
  \item Pink.
  \item White.
  \item Beige (米色).
\end{choices}
\question{16. What furniture are the speakers planning to replace?}
\begin{choices}(3)
  \item The sofas.
  \item The bookshelf.
  \item The cupboards.
\end{choices}
\question{17. What will the speakers do next?}
\begin{choices}(3)
  \item Get the exhibition's details.
  \item Set off for the gallery.
  \item Buy a painting online
\end{choices}
\par 听第 10 段材料, 回答第 18 至 20 题.\par
\question{18. What is the woman trying to tell us about eye health?}
\begin{choices}(3)
  \item Some medical developments.
  \item Some healthy habits.
  \item Some false beliefs.
\end{choices}
\question{19. What do studies about wearing the wrong glasses show?}
\begin{choices}(3)
  \item It weakens adults' eyes.
  \item It is dangerous for children.
  \item It changes the eye structure.
\end{choices}
\question{20. What do we learn about eye replacement from the talk?}
\begin{choices}(3)
  \item It is not difficult.
  \item It can never be done.
  \item It can partly be done
\end{choices}
\parttitle{第二部分 阅读理解（共两节，满分50分）}
\sectitle{第一节 (共 15 小题；每小题 2.5 分，满分 37.5 分)}
\par 阅读下列短文，从每题所给的 A、B、C 和 D 四个选项中，选出最佳答案。\par
\passagelabel{A}
\par Journal of Cultural Heritage, a journal of science and technology, combines diverse academic disciplines for studying problems about the awareness and conservation of cultural heritage in a wide range. The main purpose of the journal is to present innovative approaches concerning all scientific aspects related to heritage science. SPECIFIC TOPICS\par
\par Analysis and conservation of heritage assets (资产): novel methods for studying the composition, dating, origin and new materials for the preservation of objects.\par
\par Conservation of built heritage: analysis of historical materials and construction techniques; novel inspection, testing and monitoring techniques.\par
\par Interaction between heritage items and the environment: the impact of climate change and risk\par
\par assessment of cultural heritage\par
\par Digital technologies for knowledge, conservation and restoration: data analysis and modelling, long-term preservation of heritage assets.\par
\par Museum conservation: technologies for the management and improvement of museum collections. REQUIREMENTS\par
\par The articles must be of great interest for a wide audience. Thus, it is announced that the number of articles dealing with case studies will be reduced in order to favor original articles. Reports on restoration activities should present a specific technical or scientific innovation.\par
\par BENEFITS TO AUTHORS\par
\par We provide several author benefits after acceptance, such as free PDFs, a broad copyright policy and special earnings as follows.\par
\par\vspace{0.4em}
\begin{center}\begin{tabular}{{|l|p{\dimexpr\linewidth-3em-4\tabcolsep}|}}
\hline
Details & Information \\
Deadline & the 15thof every month \\
Payment & \$60 per thousand words \\
Early submission bonus & a \$15 bonus if submitted before the 10thof every month \\
\hline
\end{tabular}\end{center}
\vspace{0.4em}
\par Please see our Guide for Authors for information on article submission. If you require further information or help, please visit our Support Center.\par
\question{21. What does Journal of Cultural Heritage aim to do?}
\begin{choices}(2)
  \item Collect technical and scientific articles.
  \item Combine kinds of academic disciplines.
  \item Offer novel methods of heritage science.
  \item Promote the growth of heritage tourism.
\end{choices}
\question{22. Which article is more likely to be included in the journal?}
\begin{choices}(2)
  \item How to survive from the economic crisis
  \item What we can do to produce digital assets
  \item Science is improving museum collections
  \item The environment is waiting for our action
\end{choices}
\question{23. How much would an author get for an accepted 2,000-word article submitted on April 5th ?}
\begin{choices}(2)
  \item \$120 .
  \item \$135 .
  \item \$165 .
  \item \$180 .
\end{choices}
\passagelabel{B}
\par Wildlife photographer Ingo Arndt discovered his first enormous ant mound (蚁丘) as a child exploring the forest with his father near their home almost 50 years ago. They were bird-watching and came around a bend in a thickly wooded area when there it was: a five-foot-tall mound, standing upward like a large stalagmite (石笋). Arndt wanted to investigate more closely, but a very particular smell suggested he rethink that. The air felt thick and sharp, discomforting his nose. "All my life," he says, "I could remember the smell".\par
\par Several years later, Arndt moved to the countryside. He continued his fascination with the mounds and their armies of tiny engineers. Outfitted with a high-resolution camera, he began photographing the mounds and sharing his imagery with researchers for a scientific perspective. It turned out the mound makers were indeed special. They were red wood ants—one of the smallest of all so-called keystone species. Keystone species such as elephants and sharks are watched closely because their behaviors affect so many aspects of the ecosystem that if they disappeared, it would struggle to adapt. Researchers also explained the smell from the mound. As the insects build a nest, they generate formic acid (蚁酸), the origin of the smell, which helps to fight against enemies\par
\par But in recent decades, the population of ants has been declining as the forests have fallen victim to logging, urbanization, and wildfires, as well as drought and higher temperatures that have become more frequent with climate change. This has led several countries across the ants' range, including Germany, to enlist them as a protected species by law.\par
\par Today Arndt's photos put on display that as social insects, these small ants not only form complex societies, but they also cohabitate with a wide variety of species, creating a mass of coexisting relationships across a variety of plant and animal species. They are also a reminder that by protecting those unnoticed creatures, we are actually protecting ourselves.\par
\question{24. How was Arndt's first encounter with the ant mound?}
\begin{choices}(2)
  \item It reminded him of a bad experience.
  \item It left him an unforgettable memory
  \item It aroused his enthusiasm for nature.
  \item It inspired him to protect the forests.
\end{choices}
\question{25. What does the underlined word "it" in Paragraph 2 refer to?}
\begin{choices}(2)
  \item The wooded area.
  \item The mound maker.
  \item The keystone species.
  \item The entire ecosystem.
\end{choices}
\question{26. What contributed to the decrease in the number of ants?}
\begin{choices}(2)
  \item Their worsened habitats.
  \item Their changing behavior.
  \item Their severe competition.
  \item Their limited population.
\end{choices}
\question{27. What is the text mainly intended to tell us?}
\begin{choices}(2)
  \item Massive lives exist on earth.
  \item Tiny species count in nature.
  \item Social creatures live in groups.
  \item Little ants are master architects.
\end{choices}
\passagelabel{C}
\par It's a classic complaint in relationships, especially romantic ones: "She said she was okay with me forgetting her birthday! Then why is she throwing dishes in the kitchen? Are the two things related?"\par
\par The answer is YES. Communication is more than words. It's how those words are said, the tone, the order, even the choice of a particular word. It's multi-dimensional, as explored by Deirdre Wilson and Dan Sperber in Meaning and Relevance. Consider the following example: Peter got angry and Mary left; Mary left and Peter got angry. Though identical in words, their meanings differ completely. The first one may have us thinking: Wow, Peter must get angry often if Mary leaves. The second sentence suggests that Peter wants more from Mary. Same word—a totally different context.\par
\par Human language is not a code (编码). True codes have a one-to-one relationship with meaning. One sound, one definition. This is what we see with animals. Wilson and Sperber explain that "coded communication works best when emitter (发出者) and receiver share exactly the same code." Not so for humans. We communicate more than the definitions of our words would suggest. This is inferential communication, and it means that we understand not only the words spoken, but the context in which they are spoken.\par
\par Irony (反语) is a great example of how powerfully we can communicate context with a few simple words. It is more than just stating the opposite. For example, when Mary says, after a boring party, "That was fun", she is neither saying literally that the party was fun nor saying "ironically" that the party was boring. Rather, she is expressing an emotion. Three words replace paragraphs, showcasing language’s efficiency.\par
\par Wilson and Sperber concluded that human language developed and became so powerful because of two unique abilities of humans: language and the power to try to interpret mental states of others. We look for context for the words we hear. And we should be very good at absorbing this context to infer meaning.\par
\question{28. What does the complaint in Paragraph 1 mainly indicate?}
\begin{choices}(2)
  \item The inconsistency between words and meaning.
  \item The importance of remembering vital moments.
  \item The necessity of straightforward communication.
  \item The unavoidability of argument in a couple's life.
\end{choices}
\question{29. What may make human language different from true codes?}
\begin{choices}(2)
  \item The context included in understanding meanings
  \item The definitions made by words in communication
  \item The coded communication between emitter and receiver
  \item The one-to-one relationship between word and meaning
\end{choices}
\question{30. How does Mary's statement "That was fun" illustrate about irony?}
\begin{choices}(2)
  \item It simplifies language to save time.
  \item It covers the true feelings of people.
  \item It needs wordy explanations to make sense
  \item It uses fewer words for complex meanings.
\end{choices}
\question{31. What does the author suggest in communication?}
\begin{choices}(2)
  \item Focusing on the opposite.
  \item Taking language literally.
  \item Reading beyond the words.
  \item Avoiding using ironic tones.
\end{choices}
\passagelabel{D}
\par Imagine waking up one morning to find your smartphone missing. You can no longer get guided to work, order your coffee, or connect with friends. For most of us, it is a nightmare (噩梦). But it should not be exactly far-fetched\par
\par Despite being more digitally connected than ever, we are lonelier and more anxious. Anxiety rates increased by 25\% globally. Depression now affects over 280 million people worldwide. All of this is due to our growing dependence on social media and smart technologies designed to offer quick relief but often at the expense of deep, meaningful engagement, which may further cause our loss of something easy to be ignored.\par
\par While these conveniences offer short-term relief, they may weaken our thinking skills. For example, GPS has undeniably transformed how we travel. But a 2020 study found that frequent GPS users struggle more with self-guided navigation. Over time, those who depended more on GPS showed a sharp decline in spacial memory and problem-solving skills. However, reducing dependence on these conveniences can lead to discomfort. In seeking comfort today, we might be trading away some discomfort.\par
\par That is the problem: the very discomfort we avoid is often what strengthens us. Viewing stress, a kind of discomfort, as a challenge rather than a threat leads to better outcomes. Research shows that mild stress can enhance adaptability. Neuroscience (神经科学) reveals that neuroplasticity—the brain's ability to form new connections—is enhanced through challenging experiences. This adaptability is crucial for developing resilience (抗逆力), a process fueled by both mental challenges and active effort.\par
\par But this doesn't mean we should reject comfort entirely; it's to balance it with intentional discomfort. Think of it as "microdosing hardship"—a series of small, manageable difficulties that keep our mental and emotional muscles strong, such as writing with a pen, or washing dishes by hand. These small acts are enough to reawaken the resilience we've buried under layers of ease. So, the next time you feel the pull of ease, pause. Ask yourself: Is this momentary comfort helping me grow or is it keeping me stuck? Choosing discomfort isn't easy, but sometimes, it's worth it.\par
\question{32. What does the underlined word "far-fetched" in Paragraph 1 probably mean?}
\begin{choices}(2)
  \item Possible to change.
  \item Certain to function.
  \item Unlikely to be settled
  \item Unable to be accepted.
\end{choices}
\question{33. What may digital conveniences cause according to the text?}
\begin{choices}(2)
  \item The fast but worthless relief.
  \item The meaningful engagement.
  \item The loss of thinking capacity.
  \item The lack of necessary comfort
\end{choices}
\question{34. Which is a real-life application of "microdosing hardship"?}
\begin{choices}(2)
  \item Taking the stairs instead of the lift.
  \item Ordering take-out food after work
  \item Exploring a new city by using GPS.
  \item Writing an email using voice-to-text.
\end{choices}
\question{35. What can be a suitable title for the text?}
\begin{choices}(2)
  \item GPS: the invisible cost of convenience
  \item Stress: the main cause of brain damage
  \item Discomfort: a hidden source of resilience
  \item Smartphones: barriers to human progress
\end{choices}
\sectitle{第二节 (共 5 小题；每小题 2.5 分，满分 12.5 分)}
\par 根据短文内容，从短文后的选项中选出能填入空白处的最佳选项。选项中有两项为多余选项。\par
\par We all know that first impressions are important and that they are lasting. That's why we work hard to make great first impressions. \blanknum{36}\par
\par First impressions might overshadow your abilities. When you are hired in an organization, a set of inferences are made about you. Making inferences is human nature. If you become known for your attention to details, you may never be seen as someone who can see the big picture. If you are the person who organizes all of the birthday celebrations, you may never be seen as management material. \blanknum{37}\par
\par First impressions can stamp you forever, often pre-determining the path of your career and how far you will be able to advance your career. These perceptions paint a picture of who people think you are. \blanknum{38} This, in large part, is why your careers are held back, why you become frustrated, dissatisfied and even want to leave the present position.\par
\par But it doesn't have to come to that if you can unlock hidden aspects of your value. Talk about what you are reading or a project you are working on at home with coworkers and your manager when appropriate. This can demonstrate other skills you possess. Embrace (接受) new things and become knowledgeable about your company and industry. \blanknum{39} Volunteer in your company. It is a good way to give yourself the opportunity for others to see you in a way that is outside of your defined role.\par
\par \blanknum{40} As people get to know you better, they will recognize your diverse interests and capabilities, potentially opening new doors for career advancement.\par
\begin{choices}(1)
  \item Here are some useful suggestions for you.
  \item Arm yourself with as much knowledge as possible.
  \item Thus, you can contribute in fresh and unexpected ways.
  \item It is often extremely difficult for people to change them.
  \item Whatever labels, they prevent people seeing your other capacities
  \item Nevertheless, the same first impressions can block our career progress
  \item By expanding your knowledge and network, you can reshape others' perceptions.
\end{choices}
\parttitle{第三部分 语言知识运用（共两节，满分 30 分）}
\sectitle{第一节 (共 15 小题；每小题 1 分，满分 15 分)}
\par 阅读下面短文，从短文后各题所给的 A、B、C 和 D 四个选项中，选出可以填入空白处的最佳选项。\par
\par The unexpected detours (迂回) in the past summer have turned an ordinary vacation into an unforgettable adventure.\par
\par It began with a simple message from my friends, "Why don't we \blanknum{41} the road this weekend?" With only a rough idea of our route for mountains and lakes, we \blanknum{42} into my car and set off into the clear night. Driving under a blanket of \blanknum{43} felt liberating; it was as if all our worries \blanknum{44} The laughter echoed (回荡) in the car as we shared funny stories in our lives.\par
\par As 45 broke over the horizon (地平线), so did our first 46 getting lost! A wrong turn led us deep into winding country roads far from where we intended to be. 47 frustrated, we quickly turned it into an opportunity for exploration instead of seeing it as a \blanknum{48} . Ne came across little towns we'd never heard of before—each with its own charm and character. Finally arriving at our \blanknum{49} destination—an extremely attractive national park—we were greeted by breathtaking views of towering mountains \blanknum{50} on glassy lakes. The \blanknum{51} was undoubtedly our hike up Pine Ridge Trail, which challenged us physically but \blanknum{52} us with views that took our breath away\par
\par This adventure \blanknum{53} not only as an escape from daily pressures but also a chance to accept uncertainty rather than \blanknum{54} it: sometimes getting lost opens doors that you didn't even know \blanknum{55}!\par
\cz{41}{hit}{fix}{cross}{block}
\cz{42}{dropped}{piled}{looked}{knocked}
\cz{43}{snow}{fog}{stars}{clouds}
\cz{44}{showed off}{turned up}{pulled in}{melted away}
\cz{45}{hope}{dream}{nightfall}{daylight}
\cz{46}{schedule}{approach}{opportunity}{challenge}
\cz{47}{Finally}{Initially}{Gradually}{Naturally}
\cz{48}{takeover}{setback}{checkout}{backup}
\cz{49}{hidden}{promised}{intended}{abandoned}
\cz{50}{standing}{rolling}{drawn}{reflected}
\cz{51}{shelter}{routine}{highlight}{adventure}
\cz{52}{impacted}{rewarded}{charged}{equipped}
\cz{53}{served}{passed}{ranked}{doubled}
\cz{54}{shy away from}{live up to}{put up with}{hold on to}
\cz{55}{mattered}{proved}{existed}{disappeared}
\sectitle{第二节 (共 10 小题; 每小题 1.5 分, 满分 15 分)}
\par 阅读下面短文，在空白处填入 1 个适当的单词或括号内单词的正确形式。\par
\par London's food scene has witnessed an inrush not only of traditional Cantonese (广东的) restaurants but also a diverse range of regional Chinese cuisines. For decades, London's Chinatown has been closely associated with classic dishes \blanknum{56} (root) in Guangdong province, while London's Chinese food has expanded far beyond that, reflecting a much \blanknum{57} (broad) cultural shift.\par
\par While Cantonese-centric offerings have laid the foundation of familiarity with Chinese culinary (烹饪的) culture, the increasing variety makes it \blanknum{58} big challenge for restaurants to meet local British tastes. To adapt to the trend, they manage to introduce new dishes and enhance existing ones \blanknum{59} (continuous) for bettering the menu. Additionally, larger culinary schools commit to teaching authentic (地道的) Chinese cooking, \blanknum{60} (strengthen) greater appreciation and mastery of the cuisine\par
\par Hotpot, a typical Chinese dish, has taken centre stage in London, described as an "adventure" \blanknum{61} locals could enjoy genuine Sichuan cuisine and unfamiliar cultural experiences. Meanwhile, hearty hand-pulled noodles of Xi'an are also making waves, driven by the growing \blanknum{62} (present) of Chinese customers and Londoners' appetite for authentic regional flavors.\par
\par Chen Xin, the chief organizer of the 2024 London Chinese Cuisine International Development Convention, \blanknum{63} (emphasize) that London was becoming a true global centre of Chinese cuisine, staying true \blanknum{64} one single mission—providing Chinese customers with a taste of home and showcasing the depth of Chinese culinary culture to international \blanknum{65} (diner).\par
\parttitle{第四部分 写作（共两节，满分 40 分）}
\sectitle{第一节 (满分 15 分)}
\par 假定你是李华，你校拟举办主题为 "Confucius in My Eyes" 的短视频征集活动，请你给校英文报写一篇作品征集启事，内容包括：\par
\par (1) 活动目的;\par
\par (2) 作品要求. 注意:\par
\par (1) 写作词数应为 80 左右;\par
\par (2) 请按如下格式在答题卡的相应位置作答.\par
\par Short Videos Wanted\par
\sectitle{第二节 (满分 25 分)}
\par 阅读下面材料，根据其内容和所给段落开头语续写两段，使之构成一篇完整的短文。\par
\par "Hey, Isaac!" Dayna said. There was a touch of awkwardness in her voice. "Feeling ready for this year's baking competition?" Dayna's initiative greetings really surprised me.\par
\par "Do I look worried?" I forced a polite smile. To be honest, I had been hoping to fix our relationship but worried about being turned down and making things worse.\par
\par Dayna and I used to be best friends. Our shared love for baking brought us together, leading us to join Mrs Cooper's baking classes. We always supported each other during the lessons, which greatly enhanced our baking skills. Last year, I cheered for her to win the championship, but we were in such an embarrassing situation just because of an argument.\par
\par "Prepare to be amazed," I said. She shot me a doubtful look. "We'll see." Today was the day of this year's competition, whose theme was "cupcakes". We entered the kitchen at the Students' Baking Club, where we attended baking classes every Saturday. Family and friends filled a makeshift seating area. After explaining the competition consisted of two rounds, Mrs. Cooper announced the official start of the competition. Immediately, all the contestants quickly dived into their work.\par
\par As I was working on my creations, I couldn't help glancing (瞥见) at Dayna now and then because I was so used to checking on her when baking. Dayna, on the other hand, was completely focused. Her hands were moving skillfully as she measured, mixed, and poured. She didn't seem to be affected by our current relationship at all, or maybe she was just better at hiding it.\par
\par Everything went very well. Soon, Mrs. Cooper announced that Dayna, two other contestants, and I were advancing to the final round. In this round, I got right into my zone. But I quickly realized I hadn't packed the chocolates that I used for decorating. It was my magic weapon for winning! Sweat rolling down my face, I didn't know what to do.\par
\par (1) 续写词数应为 150 左右;\par
\par (2) 请按如下格式在答题卡的相应位置作答.\par
\writingpara{Paragraph 1: Dayna must have sensed my panic.}
\writinglines{8}
\writingpara{Paragraph 2: Thanks to Dayna's chocolates, I won!}
\writinglines{8}
\end{document}
